\abstract{РЕФЕРАТ}

Объем работы равен \formbytotal{lastpage}{страниц}{е}{ам}{ам}. Работа содержит \formbytotal{figurecnt}{иллюстраци}{ю}{и}{й}, \formbytotal{tablecnt}{таблиц}{у}{ы}{}, \arabic{bibcount} библиографических источников и \formbytotal{числоПлакатов}{лист}{}{а}{ов} графического материала. Количество приложений~-- 2. Графический материал представлен в приложении А. Фрагменты исходного кода представлены в приложении Б.

Перечень ключевых слов: навигация, замкнутое пространство, эвакуационный маршрут, алгоритм Дейкстры, динамическая маршрутизация, навигационный граф, исключение узлов, межэтажный переход, чрезвычайная ситуация, кратчайший путь, QR-навигация, расписание, .NET MAUI, SkiaSharp, MVVM.

Объектом разработки является система построения эвакуационных маршрутов в реальном времени для замкнутых пространств, реализованная в виде кроссплатформенного приложения на платформе .NET MAUI.

Целью выпускной квалификационной работы является программная реализация и тестирование алгоритмов динамической маршрутизации, обеспечивающих построение оптимальных эвакуационных маршрутов с учётом динамически изменяющихся препятствий и режима чрезвычайной ситуации.

В ходе работы был проведён анализ алгоритмов поиска кратчайших путей на графах, разработаны структуры данных навигационного графа и механизмы перестроения маршрута в обход заблокированных узлов, реализованы алгоритм Дейкстры и алгоритм поиска ближайшего эвакуационного выхода.


Разработанная система позволяет формировать эвакуационные маршруты в реальном времени и служит основой для дальнейшего развития навигационных систем в замкнутых пространствах.

\selectlanguage{english}
\abstract{ABSTRACT}

The volume of work is \formbytotal{lastpage}{page}{}{s}{s}. The work contains \formbytotal{figurecnt}{illustration}{}{s}{s}, \formbytotal{tablecnt}{table}{}{s}{s}, \arabic{bibcount} bibliographic sources and \formbytotal{числоПлакатов}{sheet}{}{s}{s} of graphic material. The number of applications is 2. The graphic material is presented in Annex~A. Source code fragments are presented in Annex~B.

List of keywords: navigation, enclosed space, evacuation route, Dijkstra’s algorithm, dynamic routing, navigation graph, node exclusion, inter-floor transition, emergency situation, shortest path, QR navigation, schedule, .NET MAUI, SkiaSharp, MVVM.

The subject of this development is a system for constructing real-time evacuation routes for enclosed spaces, implemented as a cross-platform application on the .NET MAUI platform.

The goal of this final qualification project is the software implementation and testing of dynamic routing algorithms that ensure the construction of optimal evacuation routes, taking into account dynamically changing obstacles and emergency conditions.

During the project, an analysis of shortest-path search algorithms on graphs was conducted; data structures for the navigation graph and mechanisms for rerouting around blocked nodes were developed; and Dijkstra’s algorithm and the nearest-exit search algorithm were implemented.

The developed system allows for the generation of evacuation routes in real time and serves as a foundation for the further development of navigation systems in enclosed spaces.
\selectlanguage{russian}
